\section{A common master and birational transport}
\label{sec:global-transport}

The obstruction in Sections~\ref{sec:incomplete-gamma}--\ref{sec:two-wall}
comes from comparing two local receivers at an intermediate chamber.  We now
avoid that comparison.  A single equivariant cobordism carries both endpoint
quantum Kirwan maps, and virtual localization compares their point rows before
either endpoint Novikov completion is taken.

\subsection{The common point row}

Let \(W\) be a smooth projective variety with a \(\Gm\)-action and let
\(Y_-\) and \(Y_+\) be smooth free GIT quotients for two extreme
linearizations.  Write \(\kappa_\pm\) for the corresponding quantum Kirwan
maps and \(\tau_{Y_\pm,-}\) for Woodward's localized graph potentials.  The
derivatives
\[
 A_\pm=D\tau_{Y_\pm,-}\,D\kappa_\pm                       \tag{8.1}
\]
are the endpoint localized gauged maps; this is the localized adiabatic
identity \cite[Equation~(68)]{WoodwardQKIII}.

Suppose that a class \(a_p\in H^*_{\Gm}(W)\) has the following restrictions:
\[
 \kappa_\pm(a_p)=[p_\pm],\qquad
 a_p|_F=0                                                     \tag{8.2}
\]
for every fixed component \(F\) occurring at an intermediate polarization.
Here \(p_-\) and \(p_+\) lie in the common birational open set.

\begin{proposition}[Support collapse]
\label{prop:support-collapse}
After Rees homogenization, the two endpoint point rows have one pullback:
\[
 \rrow_{Y_-,p_-}A_- =\rrow_{Y_+,p_+}A_+.                    \tag{8.3}
\]
The identity is coefficientwise in all remaining Novikov and bulk variables.
\end{proposition}

\begin{proof}
Retain rotation of the parametrized graph curve, with equivariant parameter
\(\hbar\).  At each ordinary input marking insert the normalized equivariant
point class of \(0\in\mathbf P^1\).  Rotation localization then puts every
input on the zero-side bubble tree.  Next choose equivariant divisor classes
\(D_1,\ldots,D_s\) whose endpoint Kirwan restrictions span the two numerical
divisor spaces, and insert Woodward's Liouville classes with parameters
\(t_j\).  On a graph fixed locus their contribution is
\[
 \exp\!\left(\sum_jt_jD_j\right)
 \prod_j x_j^{D_j\cdot d_+},
 \qquad x_j=\exp(t_j\hbar),                                  \tag{8.4}
\]
where \(d_+\) is the degree on the infinity side
\cite[Lemma~9.3]{WoodwardQKIII}.  Extracting the trivial character in all
\(x_j\) forces \(d_+=0\).  The infinity factor is therefore the unstable
identity, while the zero factor is exactly the localized graph potential in
(8.1).  Using the full divisor lattice is important: a single wall-nef
divisor can vanish on an extremal ray and would not detect \(d_+\).

Every nonendpoint polarization-fixed graph has its principal component in a
fixed quotient.  Moreover every marking and node lands in the corresponding
fixed locus \cite[Proposition~3.15(c), Lemma~3.17 and
Proposition~3.18]{GonzalezWoodward}.  Its insertion therefore factors through
one of the restrictions in (8.2) and vanishes before division by the virtual
normal Euler class.  Only the two endpoints survive the virtual Kalkman
identity.  Their signs and node factors agree with Woodward's localized
adiabatic identity, giving (8.3).  Polarizing in the ordinary inputs gives
the entire row rather than only the scalar potential.
\end{proof}

For a birational map between smooth projective varieties, the required class
exists globally.  W{\l}odarczyk's smooth birational cobordism admits a smooth
projective equivariant completion whose source and sink quotients are the two
endpoints \cite[Proposition~2(B')]{Wlodarczyk}.  Choose \(p\) in the common
birational open.  There the cobordism is a trivial orbit cylinder.  The
equivariant Poincar\'e dual of the closure of that orbit restricts to the two
point classes and misses every intermediate fixed component.  It is the
class \(a_p\) in (8.2).

\subsection{One coefficient field}

The two maps in (8.1) initially belong to opposite Novikov completions.  The
next theorem compares them before completing in the exceptional directions.
Fix a finite ample-energy and bulk Artin quotient of the equivariant Novikov
ring.  For a primitive affine gauge direction \(\delta\), all ordinary curve
degrees and fixed graph types are then finite; only the integer coordinate
along \(\delta\) can remain unbounded.

\begin{theorem}[Complete-neutral localization]
\label{thm:complete-neutral}
Suppose
\[
 c_1^{\Gm}(TW)\cdot\delta=0.                                 \tag{8.5}
\]
For every fixed ordinary curve class, the complete sum of rotation-fixed
terms in the localized gauged potential is an Artin-valued balanced
Mellin--Barnes kernel.  The two extreme polarization series are its two
oriented residue expansions.  The assertion is preserved by arbitrary input
and bulk derivatives and by the Rees substitution
\[
 Q^d\longmapsto z^{c_1^{\Gm}(TW)\cdot d}Q^d.                 \tag{8.6}
\]
If (8.5) fails, each homogeneous Artin coefficient is Laurent-finite in the
\(\delta\)-direction.  Consequently the two maps \(A_-\) and \(A_+\) coexist
over one meromorphic differential field.
\end{theorem}

\begin{proof}
Woodward's classification of the rotation-fixed graph locus is exhaustive
\cite[Corollary~9.10]{WoodwardQKIII}.  His virtual-normal splitting gives
\[
 \operatorname{Eul}(N_\pm)=
 \operatorname{Eul}\!\left(
  (R\pi_*\operatorname{ev}^*T(W/\Gm))_{\mathrm{mov}}
 \right)(\mp\zeta)(\mp\zeta-\psi).                           \tag{8.7}
\]
The last two factors smooth the node and move the attaching point.  They are
independent of the affine gauge degree, so gluing neither removes one end of
the degree lattice nor creates a degree-dependent pole.

Pass to a finite cover which clears stabilizer denominators and apply the
splitting principle.  A virtual line in the moving index of the principal
weighted component has nilpotent Chern root \(\alpha_a\) and degree
\[
 n_a(k)=h_ak+s_a.                                             \tag{8.8}
\]
For every integer \(n\), its inverse index Euler factor is the single
meromorphic expression
\[
 \frac{1}{\prod_{m=0}^{n}(\alpha_a+m\zeta)}
 =\zeta^{-n-1}
   \frac{\Gamma(\alpha_a/\zeta)}
        {\Gamma(\alpha_a/\zeta+n+1)}.                         \tag{8.9}
\]
For \(n<0\), the left side denotes the Euler class of \(H^1\); for
\(n\geq0\), it denotes the inverse Euler class of \(H^0\).  Thus the two
degree rays are not assigned unrelated Gamma products.

The signed sum of the moving slopes is
\[
 \sum_a\epsilon_a h_a=c_1^{\Gm}(TW)\cdot\delta.              \tag{8.10}
\]
The gauge Lie algebra has slope zero.  Moving modes on attached fixed bubbles
have fixed rank at the chosen ordinary degree.  Their Euler factors are
finite products of linear functions of \(k\), possibly inverted.  Each
inverse factor is an adjacent Gamma ratio
\[
 (hk+a)^{-1}=\frac{\Gamma(hk+a)}{\Gamma(hk+a+1)},              \tag{8.11}
\]
so it has zero signed slope and does not alter (8.10).  Equivariant inputs
are polynomial in \(k\); Chern roots and psi classes are nilpotent at Artin
level; Liouville classes supply the Fourier exponential.

Under (8.5), these factors satisfy the balanced contour hypotheses of
Aleshkin--Liu's Mellin--Barnes wall-crossing theorem
\cite[Definition~5.18 and Theorem~5.21]{AleshkinLiu}.  The present use is not
a direct quotation of their linear GLSM theorem: equations (8.7)--(8.11)
reduce each nonlinear fixed-graph coefficient to their balanced kernel.
The residue orientation is fixed by virtual localization.  For one moving
character,
\[
 \operatorname*{Res}_{h\sigma+\alpha=-m}
 \Gamma(h\sigma+\alpha)\,d\sigma
 =\frac{(-1)^m}{h\,m!}.                                      \tag{8.12}
\]
The factorial is the moving-mode Euler factor, \(1/h\) is the
stabilizer/Jacobian factor, and the sign is the complex virtual-normal
orientation.  Products of (8.12), followed by the finite nilpotent
derivatives, reproduce every fixed term with its multiplicity.

One must sum the gluing types before moving the contour.  A pole in (8.11)
can be the boundary presentation of another fixed graph.  Woodward's cutting
and gluing identities and the exhaustive fixed-locus union show that the
complete finite graph sum contains each such boundary term once.  Contour
deformation therefore produces the two complete chamber expansions, not
merely two sequences satisfying the same recurrence.

If (8.5) fails, the virtual-dimension equation fixes \(k\) in each homogeneous
coefficient, hence the coefficient is Laurent-finite.  Taking a basis of the
degree lattice combines the Laurent-finite and balanced directions into one
field.  Differentiation and (8.6) only add polynomial factors or nilpotent
Gamma shifts.  All operations commute with quotient maps between Artin
truncations, so the inverse limit never evaluates a formal Novikov variable
at a nonzero value.
\end{proof}

\subsection{The finite packet}

Formal smooth-endpoint quantum Kirwan surjectivity is elementary in this
setting.  Modulo the positive-energy ideal, \(D\kappa_\pm\) is the derivative
of the classical Kirwan map and is surjective.  Choose a linear section of
the reduction.  Its arbitrary lift has composite \(1+N\), with \(N\) in the
augmentation ideal, and the convergent Neumann series
\(\sum_{j\geq0}(-N)^j\) gives a right inverse.  Since
\(D\tau_{Y_\pm,-}\) is an invertible fundamental solution, the maps (8.1)
are surjective as well.

The Rees factor in (8.6) is needed to compare formal monodromy.  Homogeneity
of \(D\kappa\) gives
\[
 \mu_YA-A\mu_W=\frac12A-\Theta_QA,                            \tag{8.13}
\]
because the acting group has dimension one and
\(\Theta_Q(Q^d)=c_1^{\Gm}(TW)\cdot d\,Q^d\).  After (8.6), the
\(\Theta_Q\)-term becomes \(z\partial_z\), so both endpoint maps intertwine
the \(z\)-connections with the same half-Tate shift.  The Euler-product term
intertwines because the virtual-dimension identity makes the quantum Kirwan
map homogeneous and its derivative a quantum-product morphism.  The target
bulk point \(\kappa(0)\) has positive energy and is formally parallel to the
large-radius point; this transport preserves the point row.

Let \(M\) be the common continued source supplied by
Theorem~\ref{thm:complete-neutral}.  It need not have finite rank.  Set
\[
 \overline M=M/(\ker A_-\cap\ker A_+).                        \tag{8.14}
\]
It injects into the direct sum of the two finite endpoint modules and is
therefore finite.  Both kernels are stable under the common connection, so
formal monodromy acts on \(\overline M\).  Primary decomposition and
surjectivity show that each endpoint \(\lambda\)-primary packet is the image
of the same shifted primary packet of \(\overline M\).

\begin{theorem}[Birational invariance of the point-row primary Boolean]
\label{thm:birational-point-primary}
Let \(Y_-\) and \(Y_+\) be smooth projective birational complex varieties.
For every formal-monodromy eigenvalue \(\lambda\),
\[
 \rrow_{Y_-}|_{\cP_\lambda(Y_-)}\ne0
 \quad\Longleftrightarrow\quad
 \rrow_{Y_+}|_{\cP_\lambda(Y_+)}\ne0.                       \tag{8.15}
\]
Here the packets are realized in the common field constructed in
Theorem~\ref{thm:complete-neutral}; no chamberwise sector identification is
part of the statement.
\end{theorem}

\begin{proof}
Use the global W{\l}odarczyk completion and its orbit-cylinder class.  By
Proposition~\ref{prop:support-collapse}, there is one row \(\rho\) on \(M\)
such that
\[
 \rrow_{Y_-}A_-=\rho=\rrow_{Y_+}A_+.                         \tag{8.16}
\]
Theorem~\ref{thm:complete-neutral} and (8.14) place both surjections and this
row in one finite differential module.  If a vector in the first endpoint
\(\lambda\)-packet is detected by the point row, lift it to the corresponding
primary packet of \(\overline M\).  Equation (8.16) shows that its image at
the second endpoint is detected by the second point row.  The reverse
implication is symmetric.
\end{proof}

Theorem~\ref{thm:birational-point-primary} is the global replacement for
Hypothesis~\ref{hyp:rank-zero-target}.  The countermodels remain relevant:
they show why the theorem cannot be obtained by composing the localized
one-wall receivers or by retaining only their unmarked Fourier boundary.
